\begin{textsample}{NEG Dim 9 – global\_south\_2025\_03 – Score -4.00 – global\_south\_2025\_03/122940986.txt}  \label{ex:f9_neg_399}
Hamas released a video of Israeli hostage Elkana Bohbot pleading for his release, warning Israeli airstrikes could kill him.
Bohbot is among 58 hostages still held.
Meanwhile, thousands in Tel Aviv protested, urging Netanyahu to secure their return.
Supporters of Israeli hostages \textbf{kidnapped} during the October 7, 2023 attack on Israel by Hamas, take part in a protest against Israeli government and its head Prime Minister Benjamin Netanyahu and to demand the release of all hostages, in Tel Aviv, Israel, March 29, 2025 REUTERS/Itai Ron ( REUTERS ) Hamas ' s armed wing on Saturday released a video showing an Israeli hostage, Elkana Bohbot, pleading for his release.
Bohbot, who was abducted from a \textbf{music} \textbf{festival} in southern Israel during Hamas ' s October 7, 2023, attack, appears in the footage calling on the Israeli government to secure his freedom.
The three-minute video, spoken in Hebrew, is the second hostage footage shared by Hamas in recent days.
The Hostages and Missing Families Forum confirmed Bohbot @ @ @ @ @ @ @ @ @ @ week in another video alongside captive Yosef Haim Ohana.
AFP stated that it was unable to verify the exact date and location of the latest footage.
Hostage fears for his life In the newly released video, Bohbot expresses his concerns over Israeli airstrikes in Gaza and the danger they pose to the hostages.
He appeals to the government for his release, saying he wishes to be reunited with his wife and son.
Since Israel resumed its offensive on March 18 after a temporary ceasefire, Hamas has warned that continued military actions could endanger the lives of the remaining hostages.
Of the 251 individuals taken captive during the October 7 attack, 58 are still held in Gaza, with the Israeli military stating that 34 of them are presumed dead.
Growing calls for a ceasefire Israel ' s renewed military campaign in Gaza has resulted in significant casualties, with nearly a dozen people killed on Saturday alone, according to the region ' s civil defense agency.
The hostilities have led to mounting calls for @ @ @ @ @ @ @ @ @ @ senior Hamas official, stated that ceasefire negotiations were gaining momentum."We hope that the coming days will bring a real breakthrough in the war situation, following intensified communications with and between mediators in recent days,"Naim said in a statement on Friday.
Protests in Israel for hostage release On Saturday evening, thousands of Israelis gathered in Tel Aviv, demanding that Prime Minister Benjamin Netanyahu ' s government reach a deal to bring the hostages home.
Some of the demonstrators included former hostages and relatives of those still in \textbf{captivity}."Soon, Israel will celebrate Passover...
I wish for us to be able to hold the seder night with the hostages, who must return so that we can truly celebrate a real \textbf{festival},"said Yair Horn, a former hostage whose brother Eitan remains in Gaza."Prime Minister... let ' s reach a deal without fighting."The upcoming Jewish \textbf{festival} of Passover, known as the"holiday of freedom,"is traditionally marked by a family @ @ @ @ @ @ @ @ @ @ Many protesters emphasised that true celebration would only be possible if all hostages were safely returned home.
The second phase aims to negotiate the release of the remaining hostages and the withdrawal of Israeli troops from Gaza.
Hamas insists any proposals must initiate this phase, while Israel has proposed extending the initial 42-day truce.

% matched lemmas: captivity, festival, kidnap, music
\end{textsample}
